\section{Proofs for Structured Expressivity}
\label{app:expressivity-proofs}

\begin{lemma}[Linear images of relative interiors]
\label{lem:linear-ri}
Let \(C\) be a nonempty convex set in a finite-dimensional vector space and
let \(L\) be linear.  Then
\begin{equation}
  L(\ri C)=\ri L(C).
  \label{eq:linear-ri}
\end{equation}
\end{lemma}

\begin{proof}
Let \(x\in\ri C\), set \(y=Lx\), and choose any \(z=Lx_z\in L(C)\).
The relative-interior extension property gives \(\eta>0\) with
\(x+\eta(x-x_z)\in C\).  Applying \(L\) gives
\(y+\eta(y-z)\in L(C)\), proving \(y\in\ri L(C)\).

Conversely, let \(y\in\ri L(C)\), choose \(x_0\in\ri C\), and set
\(z=Lx_0\).  The extension property in \(L(C)\) gives \(\eta>0\) and
\(x_1\in C\) satisfying \(Lx_1=y+\eta(y-z)\).  Then
\[
  x=\frac{x_1+\eta x_0}{1+\eta}
\]
belongs to \(\ri C\) and satisfies \(Lx=y\).
\end{proof}

\begin{proof}[Proof of \cref{thm:diagonal-expressivity}]
Let \(L:\R^d\to\Sym^m\) be \(L(p)=\sum_jp_jG_j\).  Since
\(\ri\R_+^d=\R_{++}^d\), \cref{lem:linear-ri} gives
\[
  L(\R_{++}^d)=\ri L(\R_+^d)=\ri\Kcal_{\mathrm{diag}}.
\]
This proves the feasibility criterion.  With \(y_j=\log p_j\), the objective
is \(\norm{y-y^0}_2\) and the constraint is
\(\sum_je^{y_j}G_j=R\).  A minimizing sequence is bounded in \(y\), hence
has a convergent subsequence whose limit remains feasible.  Thus the infimum
is attained.  Linear independence of the \(G_j\) gives uniqueness of the
coefficients.

The cone \(\Kcal_{\mathrm{diag}}\) is finitely generated and closed.  If
\(R\) lies outside it, closed-cone separation gives symmetric \(Y\) with
\(\tr(YR)<0\) and \(\tr(YG_j)\ge0\) for all \(j\).  If \(R\) lies on its
relative boundary, apply the supporting-hyperplane theorem in
\(S=\operatorname{span}\Kcal_{\mathrm{diag}}\).  It gives a linear
functional on \(S\), nonzero on \(S\), that vanishes at \(R\) and is
nonnegative on the cone.  Represent it by a symmetric \(Y\).  Then
\(\tr(YR)=0\) and \(\tr(YG_j)\ge0\) for every \(j\).  At least one of these
generator pairings is strictly positive: otherwise the functional would
vanish on the span of all generators, namely on \(S\).  This proves the
stated facial certificate and excludes ambient annihilators of the cone.
\end{proof}

\begin{proof}[Proof of \cref{thm:block-expressivity}]
Let \(C=\prod_b\PSD^{(b)}\) and
\[
  L((P_b)_b)=\sum_bA_b^\top P_bA_b.
\]
The relative interior of \(C\) is \(\prod_b\SPD^{(b)}\).  Therefore
\cref{lem:linear-ri} gives
\[
  L\!\left(\prod_b\SPD^{(b)}\right)
  =\ri L(C)=\ri\Kcal_{\mathrm{block}}.
\]
This proves strict block feasibility.  If a feasible sequence has bounded
objective, the generalized eigenvalues of each
\((P_b^{(n)},P_b^0)\) stay in a common compact interval
\([e^{-C_0},e^{C_0}]\).  A subsequence converges to positive definite blocks,
and the linear equality is closed, proving attainment.

For each block, the image cone
\[
  C_b=\{A_b^\top P_bA_b:P_b\succeq0\}
\]
is exactly the cone of PSD matrices whose range is contained in
\(\range(A_b^\top)\), and is therefore closed.  To see that their finite sum
is closed, let \(X_n=\sum_bX_{b,n}\) converge with \(X_{b,n}\in C_b\).
Every summand is PSD, so
\(0\le\tr(X_{b,n})\le\tr(X_n)\).  The summands are consequently bounded;
passing to common subsequences gives \(X_{b,n}\to X_b\in C_b\) and
\(X=\sum_bX_b\).  Hence
\(\Kcal_{\mathrm{block}}=\sum_bC_b\) is closed.

If \(R\notin\Kcal_{\mathrm{block}}\), separation from the closed
cone gives symmetric \(Y\) with \(\tr(YR)<0\) and
\[
  \tr\!\left(Y\sum_bA_b^\top P_bA_b\right)\ge0
  \quad\text{for all }P_b\succeq0.
\]
Blockwise duality of the PSD cone makes this equivalent to
\(A_bYA_b^\top\succeq0\) for every \(b\).  If \(R\) lies on the relative
boundary, take a supporting functional that is nonzero on
\(S=\operatorname{span}\Kcal_{\mathrm{block}}\).  It gives the same
blockwise condition with \(\tr(YR)=0\).  Some
\(A_bYA_b^\top\) must be nonzero: if every one vanished, the functional
would vanish on every image \(A_b^\top P_bA_b\), hence on \(S\).
\end{proof}
